% !PRINTED 219
\weberpage{219}

The roots of $f$ correspond to the six vertices of the octahedron or
the six faces of the cube; the roots of $W$ to the eight faces of the octahedron
or vertices of the cube, and the roots of $K$, both for the octahedron
and for the cube, to the 12 edges (cf. \S.\,48).

If one represents the forms $W$, $K$ thus:
\[
W = (x_1^4 + x_2^4)^2 + 12\,x_1^4x_2^4, \quad K = (x_1^4 + x_2^4)^3 - 36\,x_1^4x_2^4(x_1^4 + x_2^4),
\]
then the identical relation subsisting between the three functions $f$, $W$, $K$
is easily derived:
\weq{15}{W^3 - K^2 = 108\,f^4.}
The octahedral group contains, as one sees, the tetrahedral group,
and arises from it by the addition of the one substitution
 $\Theta(x) = ix$. The fundamental forms of the octahedron are therefore
at the same time invariant forms of the tetrahedron, and in fact
the forms $f$ of the tetrahedron and octahedron agree with one another
 exactly, and we have $W = \Phi_1\Phi_2$,
\[
K = \tfrac{1}{2}(\Phi_1^3 + \Phi_2^3), \quad (\S.\,56).
\]
Written with the determinant 1, the two generating
 substitutions of the octahedral group read as follows:
\[
\Theta = \begin{pmatrix} \sqrt{i}, & 0 \\ 0, & \dfrac{1}{\sqrt{i}} \end{pmatrix}, \quad
\chi = \begin{pmatrix} \dfrac{1}{\sqrt{2i}}, & \dfrac{1}{\sqrt{2i}} \\[6pt] \dfrac{1}{i\sqrt{2i}}, & \dfrac{-1}{i\sqrt{2i}} \end{pmatrix},
\]
in which $\sqrt{2i} = 1 + i$ is to be set.

Under the substitution $\Theta$ the form $f(x_1, x_2)$ now \gesperrt{changes}
its \gesperrt{sign}, as formula (3) immediately shows. By the
application of $\chi$ the linear factors of $f$ are altered as
 follows:
\[
\begin{gathered}
x_1, \quad x_2, \quad x_1 + x_2, \quad x_1 - x_2, \quad x_1 + ix_2, \quad x_1 - ix_2 \\
\frac{x_1 + x_2}{\sqrt{2i}}, \quad \frac{x_1 - x_2}{i\sqrt{2i}}, \quad -i(x_1 + ix_2), \quad (x_1 - ix_2), \quad \frac{2x_1}{\sqrt{2i}}, \quad \frac{2x_2}{\sqrt{2i}},
\end{gathered}
\]
and from this it follows that $f(x_1, x_2)$ \gesperrt{remains unchanged} under
application of the substitution $\chi$, and thereby the
changes of the form $f$ under all the other octahedral substitutions
 are determined at the same time.

The \gesperrt{Hessian} covariant $W$ of $f$ remains unchanged
when $f$ is changed into $-f$, and consequently $W$ remains under the
%CONT>>
